% GENERATED FILE. Edit canonical lessons, then run npm run generate:books.
% canonical-source-hash: fnv1a64:ccc4a3e569c3abd8
% canonical-lessons: ZH-C489-tiaocao, ZH-C489-zhigong, ZH-W489-zhi-post, ZH-C489-quanzhi, ZH-C489-zhiwei, ZH-C489-shengzhi

\chapter{On the Staff}
\label{ch:zh-on-the-staff}
\hlchaptermodality{489}

\begin{chapteropening}
\textbf{By the end of this chapter:} \emph{I can talk about changing employers, the staff, full-time jobs, posts and promotion.}

Complete the last lesson of chapter 489: I can talk about changing employers, the staff, full-time jobs, posts and promotion.
\end{chapteropening}

\section[tiàocáo]{tiàocáo — to move to another employer; to change jobs}

\label{lesson:ZH-C489-tiaocao}



\begin{quote}
Before the new one: say the Mandarin for to translate, then the Mandarin for to earn money.
\end{quote}

\subsection*{You'll want to know: tiàocáo}
\emph{tiàocáo} — \textquotedblleft{}to move to another employer; to change jobs\textquotedblright{}.

\begin{quote}
\emph{gōngzuò yālì tài dà suǒyǐ tā tiàocáo le} — The pressure at work was too much, so he moved to another firm.
\end{quote}

Its characters come later; for now it is a word for the ear.

\begin{cousinweb}
\emph{tiàocáo} is built from \emph{tiào} 'to jump' and \emph{cáo} 'a trough', both inherited from Middle Chinese, \emph{dew} and \emph{dzaw} (Guangyun, in Baxter's transcription), which glosses the first 'to leap' and the second 'a horse's trough'. A horse leaving one trough for the next is the picture usually given, but the dictionary does not give the sense.
\end{cousinweb}

\begin{grammarlens}[title={What you've built}]
One more word for saying how it was.
\end{grammarlens}

\subsection*{Your turn}
\begin{itemize}
\raggedright
  \item \emph{Say it:} \emph{tiàocáo}
  \item \emph{Say it:} \emph{tiàocáo}, once more
  \item \emph{Say it:} \emph{gōngzuò yālì tài dà suǒyǐ tā tiàocáo le} — The pressure at work was too much, so he moved to another firm.
  \item \emph{Recall:} say the Mandarin for to translate, then the Mandarin for to earn money, then say \emph{tiàocáo} again
\end{itemize}

\begin{tcolorbox}[breakable,colback=teal!4,colframe=teal!35!black,title={Before you move on}]
What does tiàocáo mean? (\textquotedblleft{}To move to another employer; to change jobs\textquotedblright{}.) Say it once more.
\end{tcolorbox}
\section[\texorpdfstring{\zh{职工} (zhígōng)}{职工 (zhígōng)}]{\zh{职工} (zhígōng) — member of staff; employee (of a firm or factory)}

\label{lesson:ZH-C489-zhigong}



\begin{quote}
Before the new one: say the Mandarin for to earn money, then the Mandarin for to move to another employer.
\end{quote}

\subsection*{You'll want to know: \zh{职工}}
\textbf{\zh{职工}} — \emph{zhígōng} — \textquotedblleft{}member of staff; employee (of a firm or factory)\textquotedblright{}.

\begin{quote}
\emph{zhè jiā gōngchǎng yǒu yìbǎi ge zhígōng} — This factory has a hundred staff.
\end{quote}

The new character, in pieces: \zh{职} — eleven strokes in two pieces: on the left an ear piece of six strokes, a bar, two verticals with two short bars between them and a stroke rising at the foot; on the right \zh{只} \emph{zhǐ}, five strokes. The next lesson puts it in your hand.

\begin{cousinweb}
\textbf{\zh{职工}} (\emph{zhígōng}) is built from \zh{职} 'a post' and \zh{工} 'work', both inherited from Middle Chinese, \emph{tsyik} and \emph{kuwng} (Guangyun, in Baxter's transcription), which glosses the first 'to be in charge of' and, after an older dictionary, 'an occupation'.
\end{cousinweb}

\begin{grammarlens}[title={What you've built}]
One more word for saying how it was.
\end{grammarlens}

\subsection*{Your turn}
\begin{itemize}
\raggedright
  \item \emph{Say it:} \emph{zhígōng}
  \item \emph{Say it:} \emph{zhígōng}, once more
  \item \emph{Say it:} \emph{zhè jiā gōngchǎng yǒu yìbǎi ge zhígōng} — This factory has a hundred staff.
  \item \emph{Recall:} say the Mandarin for to earn money, then the Mandarin for to move to another employer, then say \emph{zhígōng} again
\end{itemize}

\begin{tcolorbox}[breakable,colback=teal!4,colframe=teal!35!black,title={Before you move on}]
What does \zh{职工} mean? (\textquotedblleft{}Member of staff; employee (of a firm or factory)\textquotedblright{}.) Say it once more.
\end{tcolorbox}
\section[\texorpdfstring{\zh{职} (zhí)}{职 (zhí)}]{\zh{职} — an ear piece, then zhǐ}

\label{lesson:ZH-W489-zhi-post}



\begin{quote}
The word first, then the character inside it.

\begin{itemize}
\raggedright
  \item \emph{Recall:} say \emph{zhígōng} — \textbf{R1}, one lesson back, before it decays
\end{itemize}

\begin{quote}
\zh{只} \zh{工}
\end{quote}

\emph{Read it:} each one aloud, with nothing said first — \emph{zhǐ}, \emph{gōng}
\end{quote}

\subsection*{Script — an ear piece, then zhǐ}
\hlblockfigure{\detokenize{figures/ZH-W489-zhi-post-filmstrip.pdf}}{How \zh{职} is written, stroke by stroke}

\begin{quote}
\zh{职}
\end{quote}

\textbf{Two pieces, left before right.} Count the pieces, not the strokes.

\textbf{The left piece, six strokes.} A bar on top. A short vertical on the left, then a long vertical on the right that runs down past everything else to the foot. Two short bars between them. Last, a stroke that rises to the right across the bottom.

\textbf{The right piece, \zh{只} \emph{zhǐ}, five strokes}, as you wrote it: \zh{口} \emph{kǒu}, then a stroke falling to the left and a dot falling to the right under it.

6 + 5 = eleven.

\subsection*{Your turn}
\begin{itemize}
\raggedright
  \item \emph{Write it:} \zh{职}, all eleven, with the model in front of you
  \item \emph{Cover:} write it again from memory, and count the pieces, not the strokes
\end{itemize}

\begin{tcolorbox}[breakable,colback=teal!4,colframe=teal!35!black,title={Before you move on}]
How many pieces does \zh{职} have? (\textbf{Two}.) What is the right piece read? (\emph{zhǐ}.) Which stroke of the left piece runs longest? (\textbf{The vertical on its right}.)

Source: \href{https://raw.githubusercontent.com/chanind/hanzi-writer-data/68d10a4b21150cae5e1ebbd223eed289cf32d90c/data/\%E8\%81\%8C.json}{Hanzi Writer Data \zh{职}, pinned stroke paths and medians}.
\end{tcolorbox}
\section[\texorpdfstring{\zh{全职} (quánzhí)}{全职 (quánzhí)}]{\zh{全职} (quánzhí) — full-time (of a job)}

\label{lesson:ZH-C489-quanzhi}



\begin{quote}
Before the new one: say the Mandarin for to move to another employer, then the Mandarin for a member of staff.
\end{quote}

\subsection*{You'll want to know: \zh{全职}}
\textbf{\zh{全职}} — \emph{quánzhí} — \textquotedblleft{}full-time (of a job)\textquotedblright{}.

\begin{quote}
\emph{quánzhí de gōngzuò bǐ dǎgōng lèi} — A full-time job is more tiring than casual work.
\end{quote}

\begin{cousinweb}
\textbf{\zh{全职}} (\emph{quánzhí}) is built from \zh{全} 'whole' and \zh{职} 'a post', both inherited from Middle Chinese, \emph{dzjwen} and \emph{tsyik} (Guangyun, in Baxter's transcription), which glosses the first 'complete'.
\end{cousinweb}

\begin{grammarlens}[title={What you've built}]
One more word for saying how it was.
\end{grammarlens}

\subsection*{Your turn}
\begin{itemize}
\raggedright
  \item \emph{Say it:} \emph{quánzhí}
  \item \emph{Say it:} \emph{quánzhí}, once more
  \item \emph{Say it:} \emph{quánzhí de gōngzuò bǐ dǎgōng lèi} — A full-time job is more tiring than casual work.
  \item \emph{Recall:} say the Mandarin for to move to another employer, then the Mandarin for a member of staff, then say \emph{quánzhí} again
\end{itemize}

\begin{tcolorbox}[breakable,colback=teal!4,colframe=teal!35!black,title={Before you move on}]
What does \zh{全职} mean? (\textquotedblleft{}Full-time (of a job)\textquotedblright{}.) Say it once more.
\end{tcolorbox}
\section[zhíwèi]{zhíwèi — a position in a firm}

\label{lesson:ZH-C489-zhiwei}



\begin{quote}
Before the new one: say the Mandarin for a member of staff, then the Mandarin for full-time (of a job).
\end{quote}

\subsection*{You'll want to know: zhíwèi}
\emph{zhíwèi} — \textquotedblleft{}a position in a firm\textquotedblright{}.

\begin{quote}
\emph{tā xiǎng yào jīnglǐ de zhíwèi} — He wants the manager's job.
\end{quote}

Its characters come later; for now it is a word for the ear.

\begin{cousinweb}
\emph{zhíwèi} is built from \zh{职} 'a post' and \emph{wèi} 'a place', both inherited from Middle Chinese, \emph{tsyik} and \emph{hwijH} (Guangyun, in Baxter's transcription), which glosses the second 'a rank in a row'.
\end{cousinweb}

\begin{grammarlens}[title={What you've built}]
One more word for saying how it was.
\end{grammarlens}

\subsection*{Your turn}
\begin{itemize}
\raggedright
  \item \emph{Say it:} \emph{zhíwèi}
  \item \emph{Say it:} \emph{zhíwèi}, once more
  \item \emph{Say it:} \emph{tā xiǎng yào jīnglǐ de zhíwèi} — He wants the manager's job.
  \item \emph{Recall:} say the Mandarin for a member of staff, then the Mandarin for full-time (of a job), then say \emph{zhíwèi} again
\end{itemize}

\begin{tcolorbox}[breakable,colback=teal!4,colframe=teal!35!black,title={Before you move on}]
What does zhíwèi mean? (\textquotedblleft{}A position in a firm\textquotedblright{}.) Say it once more.
\end{tcolorbox}
\section[shēngzhí]{shēngzhí — to get promoted}

\label{lesson:ZH-C489-shengzhi}



\begin{quote}
Before the new one: say the Mandarin for full-time (of a job), then the Mandarin for a position in a firm.
\end{quote}

\subsection*{You'll want to know: shēngzhí}
\emph{shēngzhí} — \textquotedblleft{}to get promoted\textquotedblright{}.

\begin{quote}
\emph{tā gōngzuò le wǔ nián cái shēngzhí} — He worked five years before he was promoted.
\end{quote}

Its characters come later; for now it is a word for the ear.

\begin{cousinweb}
\emph{shēngzhí} is built from \emph{shēng} 'to rise' and \zh{职} 'a post', both inherited from Middle Chinese, \emph{sying} and \emph{tsyik} (Guangyun, in Baxter's transcription), which glosses the first as a measure of grain and 'to complete'; the sense 'to rise' is not given there.
\end{cousinweb}

\begin{grammarlens}[title={What you've built}]
One more word for saying how it was.
\end{grammarlens}

\subsection*{Your turn}
\begin{itemize}
\raggedright
  \item \emph{Say it:} \emph{shēngzhí}
  \item \emph{Say it:} \emph{shēngzhí}, once more
  \item \emph{Say it:} \emph{tā gōngzuò le wǔ nián cái shēngzhí} — He worked five years before he was promoted.
  \item \emph{Recall:} say the Mandarin for full-time (of a job), then the Mandarin for a position in a firm, then say \emph{shēngzhí} again
\end{itemize}

\begin{tcolorbox}[breakable,colback=teal!4,colframe=teal!35!black,title={Before you move on}]
What does shēngzhí mean? (\textquotedblleft{}To get promoted\textquotedblright{}.) Say it once more.
\end{tcolorbox}
